% 2010 年湖南卷文科第 13 题：体积 20 cm^3 的几何体的三视图。
% 来源 PDF：普通高考/2010/2010湖南文.pdf，SHA-256 38611c33c9da2139ad6cd4117513ee144b09a6860180eb976c446b4858ada122。
% 原页：第 1 页（A3 三栏）第 2 栏下部，书页码 603。
% 重绘依据：三个直角三角形——正视图直角边 5 与 h（直角在左下），
%   侧视图直角边 6 与 h（直角在左下），俯视图直角边 5、6（直角在左上，原页未标数字）；
%   h 按原页比例画出，仅示意，不表示答案。
\begin{tikzpicture}[scale=0.62, semithick]
  % 正视图：宽 5，高 h
  \draw (0,0) -- (5,0) -- (0,4) -- cycle;
  \draw[<->] (-0.7,0) -- (-0.7,4);
  \draw (-0.85,0) -- (-0.45,0);
  \draw (-0.85,4) -- (-0.45,4);
  \node[left] at (-0.95,2) {$h$};
  \draw[<->] (0,-0.7) -- (5,-0.7);
  \draw (0,-0.85) -- (0,-0.45);
  \draw (5,-0.85) -- (5,-0.45);
  \node[below] at (2.5,-0.65) {$5$};
  \node[below] at (2.5,-1.6) {正视图};

  % 侧视图：宽 6，高 h
  \draw (7.6,0) -- (13.6,0) -- (7.6,4) -- cycle;
  \draw[<->] (7.6,-0.7) -- (13.6,-0.7);
  \draw (7.6,-0.85) -- (7.6,-0.45);
  \draw (13.6,-0.85) -- (13.6,-0.45);
  \node[below] at (10.6,-0.65) {$6$};
  \node[below] at (10.6,-1.6) {侧视图};

  % 俯视图：宽 5，高 6，直角在左上
  \draw (0,-4.4) -- (5,-4.4) -- (0,-10.4) -- cycle;
  \node[below] at (2.5,-10.7) {俯视图};
\end{tikzpicture}
